\clearpage
\subsection{商家用例}

\srsfigurelarge{figures/srs/12-merchant-use-case-v2.png}{商家用例图}{UML 用例图}{fig:merchant-use-case}

\begin{srsusecase}{8}{管理店铺商品}{商家}{P0}{uc:merchant-store-product}
\ucrow{用例分解}{8.1 管理店铺信息；8.2 管理商品分类；8.3 管理商品库存。}
\ucrow{前置条件}{用户已通过商家角色审核；涉及店铺经营操作时店铺已经完成相应审核。}
\ucrow{基本流程}{1）商家角色审核通过后提交并维护店铺基本资料，查看店铺审核状态；2）店铺审核通过后，可在营业、休息和临时闭店之间切换营业状态；3）查看并维护本人店铺的商品分类；4）新增、修改、上下架商品，维护价格、图片和标签；5）维护库存和限购数量。}
\ucrow{异常流程}{未审核商家、未审核店铺、非店铺所有者、非法营业状态、价格小于等于零、库存小于零、限购非正整数，或删除已被有效订单引用的分类/商品时拒绝操作。}
\ucrow{说明}{商品被历史订单引用后使用逻辑下架而非物理删除；库存预占和恢复通过统一业务服务完成；非营业店铺不能创建新订单。}
\ucrow{验收标准}{商家 A 不得修改商家 B 的店铺资源；下架、售罄和超限购商品不可下单；并发订单不得使库存为负。}
\end{srsusecase}

\begin{srsusecase}{9}{处理商家订单}{商家}{P0}{uc:merchant-order}
\ucrow{用例分解}{9.1 查看本店订单；9.2 处理接单请求；9.3 确认订单出餐。}
\ucrow{前置条件}{商家已通过审核并拥有当前订单所属店铺。}
\ucrow{基本流程}{1）商家按状态查看本店订单；2）对待接单订单选择接受或拒绝；3）接受后订单进入制作中；4）制作完成后确认出餐；5）系统创建待骑手接单的配送任务并向相关用户生成通知。}
\ucrow{异常流程}{非本店订单、已取消订单、重复接单/拒单/出餐或跳过合法状态时拒绝操作。}
\ucrow{验收标准}{接单、拒单和出餐严格遵循订单状态规则；重复出餐不会创建第二条活动配送任务；商家只能操作本店订单。}
\end{srsusecase}

\begin{srsusecase}{10}{管理评价与经营}{商家}{P1}{uc:merchant-operation}
\ucrow{用例分解}{10.1 回复顾客评价；10.2 查看经营数据。}
\ucrow{基本流程}{1）商家查看本店可展示评价，并对未回复评价提交一次回复；2）按时间查看已完成订单量、营业额、客单价、热销商品和评分趋势；3）经营数据以数字卡片、趋势或排名等形式展示。}
\ucrow{扩展用例}{“获取 AI 经营摘要”作为“10.2 查看经营数据”的扩展能力，可在已有统计结果基础上生成解释性摘要，不改变订单、评价或原始统计数据。}
\ucrow{异常流程}{商家不得回复其他店铺评价，不得修改顾客评分和原文；未授权用户不得查看店铺私有经营数据；AI 服务失败时仍展示原始经营统计。}
\ucrow{验收标准}{取消订单不计入营业额；回复后顾客原评价保持不变；AI 摘要与基础统计分离且失败不影响看板使用。}
\end{srsusecase}
